\chapter{Conclusioni}
La pipeline di segmentazione di tumori cerebrali in immagini MRI con tecniche classiche di \textit{Image Processing} per il confronto tra approccio monomodale e multimodale ha mostrato le migliori prestazioni medie nella configurazione \textit{Fusion3}.

Nel complesso, è stato valutato come l'utilizzo di tecniche tradizionali di segmentazione possano fornire risultati significativi, nonostante l'adozione predominante di tecniche di \textit{Deep Learning} e Reti Neurali Convoluzionali (CNN) come stato dell'arte attuale.

\section{Limiti Strutturali}
L’analisi dei casi migliori ha mostrato che la pipeline funziona bene quando la lesione presenta una morfologia compatta e un contrasto netto rispetto al tessuto circostante. In tali condizioni, la combinazione delle diverse modalità consente di ottenere contorni aderenti alla ground truth e di limitare la presenza di falsi positivi.

Al contrario, i casi peggiori hanno messo in evidenza alcune fragilità del metodo. In particolare, la multimodalità pur essendo mediamente vantaggiosa non garantisce automaticamente un miglioramento in ogni paziente, poichè se una delle sequenze contribuisce in modo non coerente alla rappresentazione finale, essa può influenzare negativamente la sogliatura iniziale, la scelta del seme e l’intera segmentazione.

\section{Sviluppi Futuri}
Tra gli sviluppi futuri, è possibile valutare strategie di fusione multimodale più avanzate rispetto alla combinazione lineare pesata, così da valorizzare il contributo delle diverse sequenze MRI.

Inoltre, sarebbe opportuno passare da una segmentazione bidimensionale basata sulle singole slice a un approccio tridimensionale sull’intero volume MRI, per sfruttare in modo più completo la continuità della lesione.

Infine, si potrebbero confrontare le tecniche adottate in questo lavoro con altri metodi classici di Image Processing, al fine di valutarne vantaggi, limiti e possibili integrazioni rispetto alla soluzione proposta. Per lo scopo, si citano le seguenti alternative.

\subsection{K-Means}
A differenza dell’approccio implementato, che segmenta la lesione attraverso fasi sequenziali, il \textit{K-Means} suddivide i pixel in un numero prefissato di classi sulla base della similarità delle intensità. Concettualmente, vi è un’analogia con la sogliatura multilivello di Otsu, poiché in entrambi i casi si ottiene una partizione dell’immagine in regioni differenti. Tuttavia, Otsu sfrutta la distribuzione dell’istogramma per determinare soglie ottimali, mentre K-Means opera iterativamente per assegnare i centroidi e aggiornare i cluster.

\subsection{Metodi Edge-Based}
Gli approcci \textit{Edge-Based} si basano sull’individuazione dei contorni nell'immagine, non sulla crescita di regioni omogenee.

Essi differiscono da quanto proposto nel lavoro presentato, in quanto la segmentazione viene costruita progressivamente a partire da una regione iniziale e una soglia, mentre gli approcci edge-based identificano direttamente le discontinuità locali di intensità che delimitano la lesione.

In maniera simile, il \textit{Marker-Controlled Watershed} è influenzato dalla struttura dei contorni. Nelle immagini tumorali in MRI i bordi della lesione sono irregolari e sfumati, rendendo le tecniche edge-based meno stabili rispetto a un approccio \textit{Region-Based} come quello adottato.

Tra gli operatori di rilevazione dei bordi più noti si citano \textit{Sobel} e \textit{Canny}. Il primo stima il gradiente locale dell’immagine mediante maschere convoluzionali orientate, mettendo in evidenza le variazioni brusche di intensità. Un possibile utilizzo sarebbe quello di evidenziare i margini principali della lesione prima della fase di rifinitura. Presenta, però, una maggiore sensibilità al rumore e tende a produrre contorni spessi e non sempre chiusi.

Il secondo è più sofisticato e migliora il processo attraverso una fase di sfocatura gaussiana, il calcolo del gradiente, la soppressione dei non-massimi e una sogliatura. Rispetto a Sobel, produce contorni più sottili e continui, ma individua prevalentemente i bordi locali e non costruisce direttamente una regione segmentata.